\begin{afterword}
\summaryhead

The post starts from Keats's lines in \textsc{Lamia}, which complain that ``cold philosophy'' will ``Empty the haunted air, and gnomed mine'' and ``Unweave a rainbow.'' One of these things, the author says, is not like the others. The haunts and gnomes were explained away: they never existed. The rainbow was explained, and it is still there.

The distinction comes from the author's ``Righting a Wrong Question.'' A belief is explained when its cause includes the thing believed, as with socks one sees on one's feet, and explained away when it does not, as with a mirage. The author thinks anti-reductionists miss this. The classic objection, that belief in reductionism is ``the mere result of the motion of molecules,'' treats every causal account of a belief as a debunking. But the belief that one wears socks is caused, through the eye and brain, by the socks.

The source of the confusion is the Mind Projection Fallacy: hearing ``there are no fundamental rainbows'' as ``there are no rainbows.'' The author admits to having worded parts of the post this way on purpose, and ends: the rainbow is still there.

\responsehead

The distinction at the center of the post is sound and clearly put. A causal account of a belief debunks it only if the chain of causes leaves out the thing believed. The reply to the ``classic objection'' works against the objection as stated, and the post is right that ``mere'' is where the objection goes wrong; an early version by J.~B.~S. Haldane (1927) uses the same word. The reply is shown for a belief formed by seeing, not for one reached by argument, but one case is enough to show that a causal account need not debunk.

The trouble is the application to Keats, which opens the post. The author reads the lines as saying that truth destroyed haunts, gnomes and rainbows alike. The poem gives partial support: the sentence the post cuts off compares unweaving the rainbow to Lamia melting ``into a shade.'' But the lines the post quotes also say that the rainbow survives. ``We know her woof, her texture; she is given / In the dull catalogue of common things.'' What ``fly'' are its ``charms.'' In those lines Keats's complaint is that the rainbow became common, not that it ceased to exist. That complaint is the one the author answers three days later in ``Joy in the Merely Real.'' Here, the reply that the rainbow ``is still there'' answers a position the quoted lines do not take. The poem's 1909 editor, Margaret Robertson, also cautions that the passage is ``Not to be taken as a serious expression of Keats's view of life,'' which is one editor's reading.

The other opponents are not identified. ``Anti-reductionists'' are said, with ``I think,'' to miss the distinction, but none is quoted. The word also covers philosophers who hold that everything, rainbows included, is physical, and who reject reduction in the sense of deriving one science's laws from another's.

Finally, the post's test on its reader, phrasings that commit the fallacy on purpose, cannot tell a reader who commits the fallacy from one who reads ``the mine de-gnomed'' as shorthand. It shows only that the shorthand is common.

In short: a sound distinction and a good answer to an old objection, applied to a poem whose quoted lines say the rainbow became common, not that it vanished, and to anti-reductionists the post does not identify.
\end{afterword}
